\section{Future Work}
\label{sec:future}

\begin{table}[!t]
\caption{Practical deployment checks for operational PV forecasting.}
\label{tab:deployment_checklist}
\centering
\scriptsize
\setlength{\tabcolsep}{3pt}
\begin{tabularx}{\columnwidth}{@{}>{\raggedright\arraybackslash}p{0.24\columnwidth}>{\raggedright\arraybackslash}X@{}}
\toprule
Layer & Required deployment check \\
\midrule
Data ingest & Verify timezone, cadence, duplicate handling, and missing-value policy before feature generation. \\
Weather input & Record whether target-time weather comes from measured data, API forecasts, numerical weather prediction, or scenarios. \\
Cloud layers & Confirm low-, mid-, and high-cloud availability or document the fallback used to synthesize missing layers. \\
Solar features & Keep latitude, timezone, solar-time convention, and clear-sky proxy consistent between training and inference. \\
Model files & Store Stage-1, Stage-2, sensor, metadata, and feature-list files together to preserve feature order. \\
Post-processing & Apply the same clipping rule used in Eq.~\eqref{eq:postprocess} and state whether metrics are pre- or post-clipping. \\
Monitoring & Track rolling MAE/RMSE, bias, ramp-event errors, and regime-specific errors after deployment. \\
Retraining & Recalibrate after seasonal drift, weather-source changes, plant changes, or sensor replacement. \\
Reporting & Report previous-day persistence, last-value persistence, and sensor-assisted diagnostics separately. \\
\bottomrule
\end{tabularx}
\end{table}


Future work should first test transferability across additional PV plants and inverter configurations. A second site would clarify whether the learned cloud-layer behavior is location-specific or more general. Another direction is uncertainty forecasting. Point forecasts support scheduling, but prediction intervals would better capture overcast or ramp-heavy periods.

Another extension is missing-data robustness. Real plant archives often contain irregular timestamps, faulty sensor blocks, or periods where inverter data are unavailable. A production pipeline should include explicit quality flags and should report separate results for clean, imputed, and fault-filtered records. The present study uses target-channel rows retained after machine-fault, hardware-error, shutdown, and invalid-sensor periods were omitted, so a future fault-aware version should not simply train on corrupted records. Instead, it should classify whether a period is a weather-driven low-power event, a sensor-quality event, or an equipment-state event, and then decide whether to forecast, impute, or exclude that interval.

The baseline study also suggests a natural next experiment. The proposed weather-input path should be tested at several explicit forecast horizons: one-step-ahead, 30-min-ahead, 1-h-ahead, and day-ahead. At very short horizons, live last-value and ramp-aware persistence baselines may be extremely strong. At longer horizons, previous-day, clear-sky-index persistence, and weather-forecast models become more appropriate. Reporting horizon-specific skill would make the operational use case sharper.

Table~\ref{tab:deployment_checklist} summarizes the engineering checks that should accompany any deployment-oriented extension. These checks are included because the workflow already supports a forecast-CSV inference path, so the next research step is to harden the data contract around timezone, weather-source type, feature availability, post-processing, monitoring, and retraining triggers, not simply to add a larger model. For a plant operator, a smaller transparent model with a reliable deployment contract may be preferable to a larger model whose inputs are not reproducible.

Finally, the forecast can be connected to microgrid scheduling and demand-response studies. In that setting, the PV forecast would serve as an input to load shifting, battery scheduling, or residential energy-management algorithms rather than being merged into the forecasting model itself. Game-theoretic or demand-response optimization could be explored in a follow-up paper, but it should not dilute the present manuscript's narrower forecasting contribution.
